% Source: https://artofproblemsolving.com/community/c3h2949413_2019_emcc_guts_round_p13_p24__exeter_math_club_competition
% Stable identity: aops-topic-2949413-q17
Write $BC=L$. Since $BP=\dfrac{L}{4}$, the right-triangle altitude theorem gives $AB^2=BP\cdot BC=\dfrac{L^2}{4}$, so $AB=\dfrac{L}{2}$ and $AC=\dfrac{\sqrt3L}{2}$. The angle bisector theorem yields $BD=\dfrac{L}{1+\sqrt3}$. This lies between $BP=\dfrac{L}{4}$ and $BM=\dfrac{L}{2}$. Thus $\dfrac{PD}{DM}=\dfrac{\dfrac{1}{1+\sqrt3}-\dfrac{1}{4}}{\dfrac{1}{2}-\dfrac{1}{1+\sqrt3}}=\boxed{\dfrac{\sqrt3}{2}}$.
\begin{asy}
usepackage("amsmath");
size(7cm);
pair B=(0,0),C=(4,0),P=(1,0),A=(1,sqrt(3)),M=(2,0),D=(4/(1+sqrt(3)),0);draw(A--B--C--cycle);draw(A--P^^A--D^^A--M);label("$A$",A,N);label("$B$",B,SW);label("$C$",C,SE);label("$P$",P,S);label("$D$",D,SE);label("$M$",M,S);
\end{asy}
